\documentclass[utf8]{beamer}
\usepackage{amsmath}
\usepackage[ngerman]{babel}
 
%Zeichnen
\usepackage{tikz}
\usetikzlibrary{calc,arrows,matrix,positioning,snakes,patterns,fadings,trees,shapes,mindmap}

\title{Learn Linux the Hard Way}
\author{Christian Hundt, Moritz Schlarb}
\institute{JGU}
\usetheme{CambridgeUS}
\beamertemplatenavigationsymbolsempty

\begin{document}

\begin{frame}
\tikzfading[name=fade outside,
           left color=transparent!0,
           right color=transparent!100,
           middle color=transparent!50]
\begin{tikzpicture}[remember picture,overlay]
  \node[yshift=-3.5cm,xshift=6.3cm] at (current page.north west) {%
    \begin{tikzpicture}
        \definecolor{unired}{rgb}{0.7568, 0, 0.176}
        \draw (0.3,2) node[right]{\includegraphics[scale=0.2]{logo}};
        \draw[fill,
              pattern color=black, 
              path fading=fade outside] (4.4,2.15) rectangle (11,1.85);
        \draw (2.6,2.65) node [right] {\begin{huge}
           \textbf{\textsf{Learn Linux }}\end{huge}};
        \draw (4.7,1.25) node [right] {\begin{Huge}\textbf{\textsf{The Hard Way}}\end{Huge}};
    \end{tikzpicture}
    
   };
 \end{tikzpicture}
 
 \begin{center}
 \vspace{28pt}
 \end{center}
\begin{center}
\begin{tikzpicture}[remember picture,overlay]
  \node[yshift=-7.7cm,xshift=5.9cm] at (current page.north west) {
\includegraphics[scale=0.9]{uni.pdf}};
\end{tikzpicture}
\end{center} 
 
\end{frame}

\section{Organisatorisches}

\begin{frame}
  \frametitle{\insertsectionhead\ -- Veranstaltung}  
  
  \begin{block}{Veranstaltungshomepage}
  \begin{center}\url{https://github.com/joguinf/learnlinuxthehardway/wiki/}\end{center}
  \end{block}
  
  \begin{block}{Mailingliste (verpflichtend)}
  \begin{center}\url{learnlinuxthehardway@lists.uni-mainz.de}\end{center}
  \end{block}
  
  \begin{block}{Archiv, RSS-Feeds, Anmeldung, Abmeldung}
  \begin{center} \url{https://lists.uni-mainz.de/sympa/info/learnlinuxthehardway} \end{center}
  \end{block}
  \begin{center}  
  \includegraphics[scale=0.2]{znurt.png}
  \end{center}
\end{frame}

\begin{frame}
  \frametitle{\insertsectionhead\ -- Scheinkriterien}
  
  \begin{block}{1. Seminarvortrag}
  Der Vortrag wird frei gehalten und dauert 60 Minuten inklusive Live-Demo der behandelten Themen. Fragen über den Stoff des Vortrags hinaus vom Betreuer und vom Publikum sollten beantwortet werden können.
  \end{block}
  
  \begin{block}{2. Schriftliche Ausarbeitung}
  Die Ausarbeitung wird in \LaTeX geschrieben, umfasst mindestens 15 Seiten, zitiert {\bf veritable} Quellen und soll im Anschluss allen Teilnehmern des Seminars frei zur Verfügung gestellt werden.
  \end{block}
  
 
 \begin{block}{3. Werkstück}
 Das Werkstück besteht aus der virtuellen Maschine, die zum Zwecke der Bewältigung der gestellten Aufgaben vom Studenten angefertigt worden ist. Auch sie ist den Seminarteilnehmern frei zur Verfügung zu stellen.
 \end{block}
 
\end{frame}
 
\begin{frame}
  \frametitle{Gentoo -- Die Metadistribution}

  \begin{columns}
    \column{.57\textwidth}
        \begin{itemize}
          \item \url{http://www.gentoo.org}
          \item Source-Distribution 
          \item Rolling Releases
          \item Modular, Customizable, Tweakable
          \item Distributionsbaukasten
          \item Hervorragende Dokumentation
          \item Versierte und hilfsbereite Community
          \item Ausgangspunkt für erste populäre Ports wie Mac Book Pro Retina, XBox, Wii, PS3, Gentoo BSD
        \end{itemize}
    \column{.43\textwidth}
      \hfill
      \begin{center}\includegraphics[scale=0.25]{Gentoo_logo.png}\end{center}
  \end{columns}

\end{frame}

\section{Themen}

\begin{frame}
  \frametitle{\insertsectionhead\ -- Basissystem}
  
\begin{columns}
\column{.45\textwidth}
\begin{block}{Basissystem A}
  \begin{itemize}
    \item Installationsmethoden
    \item Wahl der Architektur
    \item Netzwerk einrichten, \texttt{sshd}
    \item Partitionierung, Wahl des Dateisystems 
    \item Entpacken des Stage3-Archivs
    \item \texttt{chroot} -- password overwrite made easy
  \end{itemize}
\end{block}

\column{.45\textwidth}
\begin{block}{Basissystem B}
  \begin{itemize}
    \item Portage-Überblick
    \item \texttt{/etc/make.conf} und \texttt{/etc/portage}
    \item USE-Flags, CFLAGS, MAKEOPTS
    \item Lokalisierung
    \item \texttt{/etc/fstab}, \texttt{/etc/conf.d}
    \item Der Kernel - manuelle Installation vs. Genkernel
  \end{itemize}
\end{block}
\end{columns}
\end{frame}

\begin{frame}
  \frametitle{\insertsectionhead\ -- Desktop \& Development}
  
\begin{columns}
\column{.45\textwidth}
\begin{block}{Desktop}
  \begin{itemize}
    \item \texttt{/etc/X11/xorg.conf}
    \item QT4 vs. GTK3 Vergleich
    \item Enlightenment DR17, Xfce, Fluxbox, Gnome 3, KDE 4
    \item Verschlüsselung des Heimatverzeichnisses
    \item Eclipse Juno, Gedit, Kate
    \item Conky, Yakuake, Guake
    \item Spiele \& 3D-Beschleunigung:
    Warsow, Dosbox, Wine
  \end{itemize}
\end{block}

\column{.45\textwidth}
\begin{block}{Development}
  \begin{itemize}
    \item Makefiles, Cmake
    \item gcc, icc, clang
    \item sun-jdk vs. openjdk, mono
    \item Swig Interface Wrapper
    \item Ebuilds schreiben
    \item Git, SVN, gitolite, cgit
    \item IRC, Mailinglisten
    \item layman \& Overlays
    \item Crosscompiling
    \vspace{0.5\baselineskip}
  \end{itemize}
\end{block}
\end{columns}
\end{frame}

\begin{frame}
  \frametitle{\insertsectionhead\ -- Server}

\begin{block}{Administration}
  \begin{itemize}
    \item SSH, SCP, SFTP, SSHFS, VPN
    \item Screen, Tmux, Rdesktop, VNC
    \item Bash, Z-Shell, iPython
    \item Cron, Syslog
  \end{itemize}
\end{block}

\begin{columns}
\column{.45\textwidth}
\begin{block}{NAS}
  \begin{itemize}
    \item RAID
    \item LUKS
    \item Samba \& Windows 7
    \item NFS
    \item rsync, Unison, Sparkleshare
  \end{itemize}
\end{block}

\column{.45\textwidth}
\begin{block}{WWW}
  \begin{itemize}
    \item Apache, PHP, MySQL
    \item Wordpress, Drupal
    \item FTPD, IRCD
    \item Ruby, Python, Postgresql
    \item virtualenv und Turbogears 2
  \end{itemize}
\end{block}
\end{columns}
\end{frame}

\begin{frame}
  \frametitle{\insertsectionhead\ -- Sonstige}
 
 
\begin{columns}
\column{.45\textwidth}
\begin{block}{Scientific Computing}
  \begin{itemize}
    \item OpenMPI
    \item OpenMP
    \item Cuda SDK
    \item BOINC
    \item NumPy, SciPy, Octave
    \item GNUPlot
    \vspace{0.5\baselineskip}
  \end{itemize}
\end{block}
\vspace{-10pt}
\begin{center}\includegraphics[scale=0.25]{larry_the_cow-full.png}\end{center}

\column{.45\textwidth}
\vspace{-10pt}
\begin{center}\includegraphics[scale=0.2]{freebsd-logo.png}\end{center}
\vspace{-10pt}
\begin{block}{Zusätzliche Themen}
  \begin{itemize}
      \item Gentoo Prefix (z.B. mit Windows als Wirtssystem)
      \item Andere Linuxdistributionen (z.B. Kernel-Kompilieren unter Debian)
      \item Andere Unixderivate z.B. FreeBSD, Darwin
  \end{itemize}
\end{block}
\end{columns}

\end{frame}

\section{Literatur}

\begin{frame}
  \frametitle{\insertsectionhead}
  \framesubtitle{}

  \begin{block}{Gentoo}
    \begin{itemize}
      \item Gentoo Handbook: \url{http://www.gentoo.org/doc/en/handbook}
      \item Gentoo Documentation Resources: \url{http://www.gentoo.org/doc/en/}
      \item Official Gentoo Wiki: \url{http://wiki.gentoo.org}
      \item Inofficial Gentoo Wiki: \url{http://en.gentoo-wiki.com}
      \item Mailinglisten-Archiv: \url{http://www.gentoo.org/main/en/lists.xml}
      \item Gentoo Forums: \url{http://forums.gentoo.org}
    \end{itemize}
  \end{block}
  \begin{block}{Linux allgemein}
    \begin{itemize}
      \item Kernel Newbies: \url{http://kernelnewbies.org}
      \item Linux from Scratch: \url{http://www.linuxfromscratch.org}
      \item nixCraft: Linux Tips, Hacks, Tutorials: \url{http://www.cyberciti.biz}
    \end{itemize}
  \end{block}
\end{frame}

\end{document}
